\chapter{AI Instructions}

\section{Policy Overview}

AI-assisted coding tools (ChatGPT, Claude, Copilot, Gemini,
Cursor, Codeium, and similar) are \textbf{permitted} throughout
this curriculum. Security research is a field where AI proficiency
is becoming essential --- learning to use these tools effectively
is part of your training.

However, the goal is \textbf{understanding}, not automation.
If you cannot explain every line of your exploit without AI
assistance, you have not completed the exercise.

\begin{warning}
  Using AI to bypass an exercise's learning objective --- for
  example, asking it to reverse-engineer a binary so you do not
  have to learn the disassembler, or to generate an entire exploit
  chain you do not understand --- defeats the purpose of the
  curriculum and will be treated as academic dishonesty.
\end{warning}

\section{Permitted Uses}

\begin{itemize}
  \item \textbf{Concept explanation}. Ask AI to clarify a vulnerability
        class, a kernel subsystem, or a network protocol.
  \item \textbf{Scaffolding code}. Generate starter code for a harness,
        a fuzzer, or a PoC skeleton --- then modify and extend it
        yourself.
  \item \textbf{Debugging}. Explain error messages, decompiler output,
        or unexpected crash behavior.
  \item \textbf{Documentation summarization}. Condense a long RFC or
        technical blog post into actionable notes.
  \item \textbf{Translation}. Translate technical content between
        English and French (both languages are used in this curriculum).
  \item \textbf{Report drafting}. Ask for help structuring a bug bounty
        report --- but fill in the technical content yourself.
\end{itemize}

\section{Prohibited Uses}

\begin{itemize}
  \item \textbf{Complete exploit generation}. Asking an AI to write
        the full exploit for an exercise, then submitting it without
        understanding.
  \item \textbf{Bypassing exercise constraints}. Using AI to circumvent
        a restriction that the exercise explicitly imposes (e.g.\
        ``do not use automated scanners'').
  \item \textbf{Automated vulnerability discovery}. Running AI-driven
        scanning tools that replace the manual analysis you are meant
        to practice.
  \item \textbf{Writing evaluation rubrics}. You are not allowed to
        use AI to generate peer-evaluation criteria or answers.
  \item \textbf{Plagiarizing without attribution}. Copying AI output
        without citing it (see Section~\ref{sec:ai-citation}).
\end{itemize}

\section{Personal Responsibility}

When you ask an AI to generate exploit code, you become responsible
for every byte of that output before using it against any target.

\begin{enumerate}
  \item \textbf{Read the generated code} line by line.
  \item \textbf{Identify any assumptions} the AI made (architecture,
        kernel version, compiler flags, network layout).
  \item \textbf{Test in isolation} before running against an exercise
        target.
  \item \textbf{Understand the failure modes}. If the generated exploit
        crashes or hangs, you must diagnose and fix it yourself.
\end{enumerate}

\hint{If an AI generates code you do not understand, ask it to
explain. If you still do not understand after the explanation,
the code is not ready to use. Switch to a simpler approach
you can fully control.}

\section{Mandatory AI Citation}
\label{sec:ai-citation}

Every file that contains AI-generated or AI-assisted code must
include a citation comment. The format is:

\begin{minted}{python}
# Generated with assistance from <tool> on <date>.
# Prompt: <one-line summary of the prompt>
# Verified by: <method used to verify the output>
\end{minted}

\subsection*{Examples}

\textbf{Python}:
\begin{minted}{python}
# Generated with assistance from Claude 3.5 Sonnet on 2025-03-14.
# Prompt: "Write a Python script that sends HTTP requests with
# modified price fields to /api/checkout/pay"
# Verified by: manual review of each request and response
\end{minted}

\textbf{C}:
\begin{minted}{c}
/* Generated with assistance from GitHub Copilot on 2025-03-14.
 * Prompt: "Implement a stack-based buffer overflow with
 * ret2libc for x86_64 Linux"
 * Verified by: compiled with -fno-stack-protector and
 * tested against local binary with ASLR disabled */
\end{minted}

\textbf{Markdown (report)}:
\begin{minted}{html}
<!--
Generated with assistance from ChatGPT-4 on 2025-03-14.
Prompt: "Structure a bug bounty report for IDOR with
sections: summary, steps, impact, mitigation"
Verified by: all technical content written manually,
AI only suggested structure
-->
\end{minted}

\subsection*{Multiple AI interactions}

If you used multiple AI tools, or iterated through several prompts,
cite each significant interaction. A group of related prompts can
be summarized:

\begin{minted}{html}
<!--
AI assistance used in this file:
- Claude 3.5 Sonnet (2025-03-14): harness skeleton generation
- ChatGPT-4 (2025-03-14): crash triage help
- Copilot (2025-03-15): inline autocomplete for loop logic
All generated code reviewed and modified manually.
-->
\end{minted}

\section{Peer Evaluation and AI}

During peer evaluation, the evaluator may:

\begin{itemize}
  \item Ask you to \textbf{explain any AI-generated code} line by line.
  \item Ask you to \textbf{reproduce the reasoning} without AI assistance.
  \item Ask to see your \textbf{prompt history} for a given exercise.
  \item Deduct points if AI citation is missing or if you cannot
        explain code that AI generated.
\end{itemize}

\begin{warning}
  Claiming AI-generated code as your own work (failure to cite)
  is plagiarism. If the evaluator suspects plagiarism, they may
  escalate to the pedagogical team for review.
\end{warning}

\section{Ethics Reminder}

\begin{itemize}
  \item AI models can produce incorrect, dangerous, or illegal code.
        You are \textbf{responsible} for verifying every output before
        execution.
  \item AI models trained on public security research may reproduce
        verbatim code from GPL-licensed projects. Check licenses
        before incorporating into your work.
  \item Do not input sensitive or proprietary code into cloud-based
        AI tools. Use local models (Llama, Mistral, DeepSeek) for
        work involving non-public targets.
  \item The security research ethics you learn here apply to AI
        usage: do not ask an AI to generate malware, ransomware,
        or code intended to harm systems outside the exercise scope.
\end{itemize}

\info{Proficiency with AI tools is a marketable skill in modern
bug bounty hunting. Use them wisely, cite them honestly, and
always maintain your own understanding of every exploit you
deploy.}
